\subsection{有限线性群的不变量：模性质}

设 K 是一个交换环，V 是一个 K-模，G 是一个作用于 V 的群。我们知道 V 的每个自同构唯一地开拓为对称代数 \(S = \mathbf{S}(V)\) 的自同构，从而 G 作用于这个代数。若 \(x \in S\) 且 \(g \in G\) ，用 \(g_{S}.x\) 表示 x 在 g 下的像。设 R 是 S 中由在 G 下不变的元素构成的子代数 \(S^{G}\) 。

假设 G 有限，V 有限型，且 K 是 noether 的。于是 S 是一个有限型 R-模，而 R 是一个有限型 K-代数（《交换代数》第 V 章，§1，第 9 号，定理 2）。假设 S 是整的，并设 N 是它的分式域。R 的分式域 L 是 N 中在 G 下不变的元素所成的集合（出处同前，命题 23 的推论），故 N 是 L 的一个 Galois 扩张。N 的每个元素都可写成 z/t，其中 \(z \in S\) 且 \(t \in R\) （出处同前，命题 23）。由《代数》第 II 章 §7 第 10 号命题 26 的推论 3，R-模 S 的秩是 {[}N : L{]}。假设 G 忠实地作用于 V。于是 N 在 L 上的 Galois 群可与 G 等同，故 {[}N : L{]} = Card G；从而

\[
\operatorname{rk} _ {\mathsf {R}} (\mathrm{S}) = [ \mathrm{N}: \mathrm{L} ] = \operatorname{Card} (\mathrm{G}).\tag{6}
\]

对任意分次代数 \(A = A_{0} \oplus A_{1} \oplus \cdots \oplus A_{n} \oplus \cdots\) ，用 \(A_{+}\) 表示理想 \(\bigoplus_{n > 0} A_{n}\) 。

定理 1. 设 K 是一个交换域，V 是 K 上一个有限维向量空间，S = S(V) 是 V 的对称代数，G 是 V 的一个有限自同构群，R 是 S 中由在 G 下不变的元素组成的分次子代数。假设 G 由伪反射生成（§2，第 1 号），且 q = Card(G) 与 K 的特征互素。则 R-模 S 具有一个由 q 个齐次元素组成的基。

\begin{enumerate}
\def\labelenumi{\alph{enumi})}
\item
  由于 \(S/(R_{+}S)\) 的每个子模在 \(R_0 = K\) 上都是自由的，只需（《代数》第 II 章，§11，第 4 号，命题 7）证明从 \(R_{+} \otimes_{R} S\) 到 \(S\) 的典范同态是单射。对任意 R-模 \(E\) ，用 \(T(E)\) 表示 R-模 \(\text{Ker}(R_{+} \otimes_{R} E \to E)\) （*换句话说，\(T(E) = \text{Tor}_{1}^{R}(R/R_{+}, E)_*\) ）。若 \(E, E'\) 是两个 R-模而 \(u\) 是从 \(E\) 到 \(E'\) 的同态，则同态 \(1 \otimes u\) （从 \(R_{+} \otimes E\) 到 \(R_{+} \otimes E'\) ）通过限制到 \(T(E)\) 上定义一个从 \(T(E)\) 到 \(T(E')\) 的同态，我们记之为 \(T(u)\) 。若 \(u'\) 是从 \(E'\) 到某个 R-模 \(E''\) 的同态，则有 \(T(u' \circ u) = T(u') \circ T(u)\) 。于是，若 G R-线性地作用于 \(E\) ，则 G 作用于 \(T(E)\) 。
\item
  群 G R-线性地作用于 S，从而也作用于 T(S)。此外，T(S) 具有一个分次 S-模的自然结构。我们首先证明，若 \(g \in G\) ，则 g 把 T(S) 的每个元素 x 变为模 \(S_{1}T(S)\) 与 x 同余的元素。只需对 g 是伪反射的情形证明这一点。此时存在 V 中一个非零向量 v，使得 \(g(x) - x \in Kv\) （对所有 \(x \in V\) ）。由于 V 生成 S，由此可知 \(g_{S}\) 平凡地作用于 S/Sv。于是，对任意 \(y \in S\) ，存在 S 中一个元素 \(h(y)\) 使得
\end{enumerate}

\[
g _ {\mathrm{S}} (y) - y = h (y) v.
\]

由于 S 是整的且 v 非零，这个元素由 y 唯一确定；立即可见 h 是 R-模 S 的一个次数为 -1 的自同态。于是 \(g_{S} - 1_{S} = m_{v} \circ h\) ，其中 \(m_{v}\) 表示 S 中比率为 v 的位似。由此，

\[
\mathrm{T} (g _ {\mathrm{S}}) - 1 _ {\mathrm{T(S)}} = \mathrm{T} (g _ {\mathrm{S}} - 1 _ {\mathrm{S}}) = \mathrm{T} (m _ {v}) \circ \mathrm{T} (h),
\]

其像含于 \(vT(S)\) ，这就证明了我们的断言。

\begin{enumerate}
\def\labelenumi{\alph{enumi})}
\setcounter{enumi}{2}
\tightlist
\item
  我们证明 \(T(S)\) 中在 G 下不变的任何元素为零。事实上，设 Q 是 R-模 S 的由下式定义的自同态
\end{enumerate}

\[
\mathrm{Q} (y) = q ^ {- 1} \sum_ {g \in G} g _ {\mathrm{S}} (y)
\]

（对所有 \(y \in S\) ）。则 \(Q(S) = R\) 。于是我们可以写 \(Q = i \circ Q'\) ，其中 \(Q'\) 是从 R-模 S 到 R-模 R 的同态，而 i 表示 R 到 S 中的典范单射。从而 \(T(Q) = T(i) \circ T(Q')\) ，且 \(T(Q') = 0\) ，因为 \(T(R) = \text{Ker}(R_{+} \otimes R \to R) = 0\) 。于是

\[
0 = \mathrm{T} (\mathrm{Q}) = q ^ {- 1} \sum_ {g \in \mathrm{G}} \mathrm{T} (g _ {\mathrm{S}}).
\]

而 \(q^{-1} \sum_{g \in G} T(g_S)\) 保持 \(T(S)\) 中在 G 下不变的元素不动。因此这些元素为零。

\begin{enumerate}
\def\labelenumi{\alph{enumi})}
\setcounter{enumi}{3}
\tightlist
\item
  假设 \(T(S) \neq 0\) 。在 \(T(S)\) 中存在一个次数最小的非零齐次元素 \(u \neq 0\) 。由 b)，u 在 G 下不变。由 c)，u = 0。这是矛盾，故 \(T(S) = 0\) 。证毕。
\end{enumerate}

注. 1) 由《代数》第 II 章 §11 第 4 号命题 7 可知，若 \((z_{1},\ldots ,z_{q})\) 是 S 的一个齐次元素族，其典范像在 \(\mathrm{S / (R_+S)}\) 中构成 \(\mathrm{S / (R_+S)}\) 在 K 上的基，则 \((z_{1},\ldots ,z_{q})\) 是 S 在 R 上的基。

\begin{enumerate}
\def\labelenumi{\arabic{enumi})}
\setcounter{enumi}{1}
\tightlist
\item
  设 \(g\) 是 \(V\) 的一个伪反射，其阶 \(n \geqslant 2\) 有限且与 \(K\) 的特征互素。由 Maschke 定理（附录，命题 2），\(V\) 可分解为 \(D \oplus H\) ，其中 \(H\) 是由 \(V\) 中在 \(g\) 下不变的元素组成的超平面，而 \(D\) 是一条直线，\(g\) 在其上通过乘以一个本原 \(n\) 次单位根作用。当 \(K = R\) 时，这仅当 \(n = 2\) 时才可能，此时 \(g\) 是一个反射；在这种情形，定理 1 所适用的群是有限 Coxeter 群。（相反，对 K = C，定理 1 适用于某些不是 Coxeter 群的群。）\(^{3}\)
\end{enumerate}

定理 2. 沿用定理 1 的假设与记号。

\begin{enumerate}
\def\labelenumi{(\roman{enumi})}
\item
  存在 \(S\) 的一个分次向量子空间，它构成 \(R_{+}S\) 在 \(S\) 中的补，且在 \(G\) 下稳定。
\item
  设 U 是这样一个补。从 \(U \otimes_{K} R\) 到 S 的典范同态是 G-模的同构，且 G 在 U（分别为 S）中的表示同构于 G 在 K（分别为 R）上的正则表示。
\end{enumerate}

事实上，对任意整数 \(n \geqslant 0\) ，K-向量空间 \(S_n\) 与 \((R_+S) \cap S_n\) 在 \(G\) 下稳定，且由 Maschke 定理（附录，命题 2）可知存在一个 G-稳定的子空间 \(U_n\) ，它构成 \((R_+S) \cap S_n\) 在 \(S_n\) 中的补。于是 \(\sum_{n \geqslant 0} U_n\) 是 \(R_+S\) 在 \(S\) 中的一个 G-稳定的补，由此得 (i)。

设 U 是 S 的一个分次向量子空间，它在 S 中构成 \(R_{+}S\) 的补且在 G 下稳定。由注 1，K-向量空间 U 的每个基也是 R-模 S 的基，从而是 S 的分式域 N 在 R 的分式域 L 上的基。于是 L-向量空间 N 可与 \(U \otimes_{K} L\) 等同。由于 U 在 G 下稳定，这一等同与 G 的作用相容。G 在 L 上的群代数 \(L[G]\) 可与代数 \(K[G] \otimes_{K} L\) 等同。L 的 Galois 扩张 N 具有正规基（《代数》第 V 章，§10，第 9 号，定理 6），这可以解释为：N（视为 \(L[G]\) -模）同构于 G 在 L 上的正则表示的模。由于 U 在 K 上有限维，由附录命题 1 可知 \(K[G]\) -模 U 同构于 G 在 K 上的正则表示。我们的断言由此立得。

